\section{Count Good Numbers}
\label{sec:rec-good-numbers}

\problemstatement{A digit string of length $n$ is called \textbf{good} if
every digit at an \textbf{even} index (0-indexed) is one of the even
digits $\{0,2,4,6,8\}$, and every digit at an \textbf{odd} index is one
of the prime digits $\{2,3,5,7\}$. Count how many good digit strings of
length $n$ exist, modulo $10^9+7$.}

\begin{patternbox}
Despite involving ``counting all possibilities,'' this is \emph{not} a
backtracking problem -- there is no need to actually generate any digit
string, because each position's choice is completely \textbf{independent}
of every other position's choice. Whenever every position in a sequence
can be filled independently from a fixed-size choice set, the total count
is simply the \textbf{product} of each position's choice count, and that
product is best computed via \textbf{fast exponentiation}, a
straight-line recursion that shrinks an exponent by half at each step.
\end{patternbox}

\subsection*{Understanding the problem carefully}

Among the $n$ positions, $\lceil n/2 \rceil$ are at even indices ($0, 2,
4, \dots$) and $\lfloor n/2 \rfloor$ are at odd indices. Each even-index
position independently has $5$ valid choices ($0,2,4,6,8$), and each
odd-index position independently has $4$ valid choices ($2,3,5,7$). By
the multiplication principle (every combination of independent choices
is a distinct valid string), the total count is
\[
  5^{\lceil n/2 \rceil} \times 4^{\lfloor n/2 \rfloor} \pmod{10^9+7}.
\]
The only remaining challenge is computing a large power modulo $10^9+7$
efficiently, since $n$ can be large enough that computing $5^{n/2}$ by
repeated multiplication, one factor at a time, would be too slow.

\begin{ideabox}[Fast Exponentiation: Halve the Exponent, Not Decrement It]
To compute $\textit{base}^{\textit{exp}}$: if $\textit{exp}=0$, the
answer is $1$ (base case). If $\textit{exp}$ is even, recursively compute
$\textit{half} = \textit{base}^{\textit{exp}/2}$, and the answer is
$\textit{half} \times \textit{half}$. If $\textit{exp}$ is odd, compute
$\textit{half} = \textit{base}^{(\textit{exp}-1)/2}$, and the answer is
$\textit{half} \times \textit{half} \times \textit{base}$.
\end{ideabox}

\subsection*{Why halving the exponent gives $O(\log n)$ instead of $O(n)$}

A naive recursive power function that decrements the exponent by $1$ at
each call would need $O(n)$ recursive calls, each doing one
multiplication. Halving the exponent instead means the recursion depth is
only $O(\log n)$ -- each level squares an already-computed result rather
than multiplying by the base value just once, exploiting the identity
$\textit{base}^{2k} = (\textit{base}^k)^2$ to skip directly from
computing one power to a power roughly twice as large, in a single
multiplication. This is the same ``shrink faster than one unit at a
time'' idea elaborated fully in Section~\ref{sec:rec-pow}.

\subsection*{Algorithm}

\begin{enumerate}
  \item $\textit{evenCount} \gets \lceil n/2 \rceil$,
        $\textit{oddCount} \gets \lfloor n/2 \rfloor$.
  \item Compute $\textit{a} \gets 5^{\textit{evenCount}} \bmod
        (10^9+7)$ via fast exponentiation.
  \item Compute $\textit{b} \gets 4^{\textit{oddCount}} \bmod
        (10^9+7)$ via fast exponentiation.
  \item Return $(\textit{a} \times \textit{b}) \bmod (10^9+7)$.
\end{enumerate}

\subsection*{Dry run}

\dryrun{Take $n=4$.}

$\textit{evenCount}=2$ (indices $0,2$), $\textit{oddCount}=2$ (indices
$1,3$). $5^2=25$, $4^2=16$. Answer $=25\times16=400$.

Checking by hand for a tiny case, $n=1$: $\textit{evenCount}=1$,
$\textit{oddCount}=0$. $5^1=5$, $4^0=1$. Answer $=5$ -- correctly, the
$5$ single-digit good numbers are $0,2,4,6,8$ (index $0$ is even, so any
even digit qualifies, and there is no odd index to satisfy).

\subsection*{C++ implementation}

\begin{lstlisting}
#include <bits/stdc++.h>
using namespace std;
const long long MOD = 1e9 + 7;

long long power(long long base, long long exp) {
    if (exp == 0) return 1;
    long long half = power(base, exp / 2) % MOD;
    long long result = (half * half) % MOD;
    if (exp % 2 == 1) result = (result * base) % MOD;
    return result;
}

int countGoodNumbers(long long n) {
    long long evenCount = (n + 1) / 2;
    long long oddCount = n / 2;

    long long a = power(5, evenCount);
    long long b = power(4, oddCount);
    return (a * b) % MOD;
}

int main() {
    cout << "Good numbers of length 4: " << countGoodNumbers(4) << endl;
    return 0;
}
\end{lstlisting}

\begin{complexitybox}
\textbf{Time Complexity.} $O(\log n)$ for each fast-exponentiation call,
giving an overall time complexity of
\[
  O(\log n).
\]
\textbf{Space Complexity.} $O(\log n)$ for the recursion stack of fast
exponentiation.
\end{complexitybox}

\begin{takeawaybox}
Whenever a counting problem's choices at every position are independent
of each other, resist the urge to enumerate -- multiply the per-position
choice counts directly, and reach for fast exponentiation the moment that
product involves a single base raised to a large power. This pattern of
recognizing independence to avoid generation entirely recurs throughout
combinatorial counting problems.
\end{takeawaybox}
